\documentclass[11pt,a4paper]{article}

\usepackage[margin=1in]{geometry}
\usepackage{booktabs}
\usepackage{multirow}
\usepackage{hyperref}
\usepackage{xcolor}
\usepackage{natbib}
\usepackage{array}

\hypersetup{
  colorlinks=true,
  linkcolor=blue!60!black,
  citecolor=blue!60!black,
  urlcolor=blue!60!black
}

\title{%
  Canon Suppression in Corpus-Loaded Assessment: \\
  A Two-Assessor Study%
}

\author{Sophie D.\ Neilson \\
  \small Independent Researcher \\
  \small ORCID: \href{https://orcid.org/0009-0001-2430-1743}{0009-0001-2430-1743}
}

\date{September 2026}

\begin{document}
\maketitle

\begin{abstract}
We report a study comparing the assessments of a corpus-loaded evaluator and a
corpus-naive observer applied to the same behavioral data. Both evaluators were
Claude sessions. The first, the Researcher, carried what we call the research
program's \emph{canon}: a role file (a session-level instruction document), a
behavioral corpus (a vocabulary of named behavioral patterns), and the
assessment vocabulary they define. It had designed the study environment
and assessed a set of agent execution traces as part of an ongoing research
program. The second, a vanilla session carrying no canon, then reviewed the
same traces and the Researcher's record. The vanilla observer flagged 7 items as significant. Of these, the
Researcher had absorbed 3 (data present but analyzed into non-significance)
and was blind to 4 (data never reached attention). Zero items were correctly
excluded. We identify two suppression mechanisms: \emph{verdict-space
limitation} (the assessment instrument has no slot for the finding) and
\emph{canon-frame limitation} (the assessor's corpus knowledge makes the
finding appear expected or non-significant). One item is explained by the
second mechanism alone: an agent produced correctly structured eliminative
reasoning while reaching a conclusion its cited evidence did not support, and
the instrument could have caught it. Whether that conclusion was in fact wrong
is ambiguous in the specimen, so we treat the item as the sharpest
illustration of the mechanism rather than as proof of a wrong root cause. We
discuss the finding as an instance of the observer-position problem in
evaluation methodology, drawing on Bourdieu's reflexive response and
distinguishing it from Garfinkel's naive-observer fix.
\end{abstract}

\section{Introduction}

When an evaluator is deeply familiar with the system under study, that
familiarity creates a structural blind spot. The evaluator's knowledge provides
a template for what correct behavior looks like, and that template can make
genuine findings appear expected, hence non-significant, before they reach the
verdict. We use \emph{canon} to refer to the evaluator's full apparatus: its
role definition, its vocabulary of named behavioral patterns, and the assessment
rubric those patterns populate.

The evaluator is not defending the system, so the blind spot comes from
expertise rather than motivation. The evaluator's knowledge makes certain data
features \emph{invisible}: they look like what correct behavior should produce,
so they do not register as findings. Expertise, which competent evaluation
requires, also suppresses the evaluator's ability to notice certain classes of
evidence. Confirmation bias of this unmotivated kind is well
documented~\citep{nickerson1998}, and the conditions under which expert
intuition can be trusted are narrower than experts assume~\citep{kahneman2009}.

We encountered this problem concretely in Run 6 of a behavioral equivalence
assay on language model agents, whose record is released with this paper
(Section~\ref{sec:data}). The assessor, the Researcher, was a Claude session
that had designed the study environment, written the evaluation instrument,
and carried the program's canon. We then ran the same traces and
the Researcher's record through a second Claude session carrying no canon. The vanilla observer flagged 7
items as significant. The Researcher had suppressed all 7.

This paper reports that study, classifies the suppression mechanisms, and frames
the result within the observer-position literature. The contribution is the
\emph{mechanism taxonomy}: two structurally distinct suppression types with
different design responses, identified in a concrete case within AI behavioral
evaluation.

\subsection{Related Work}

Two literatures bear on the finding. In the study of expert judgment,
\citet{kahneman2009} set out the conditions under which
expert intuition is trustworthy (a regular environment and the opportunity to
learn it), and \citet{nickerson1998} reviews confirmation bias as a
phenomenon that operates without a motive to confirm. In language-model
evaluation, \citet{zheng2023} document position, verbosity, and
self-enhancement biases in model judges, and \citet{panickssery2024} show that model evaluators recognize and favor
their own generations. Both assessors in this study were Claude sessions, so
the second literature applies to both of them; the vanilla observer's flags are
not exempt from judge bias, a point the limitations return to. Neither
literature, to our reading, separates the two suppression mechanisms described
here. We did not conduct a systematic search.

\section{Study Design}

\subsection{Specimen}

The specimen was Run~6 of a behavioral equivalence assay: three agent execution
traces produced under different loading conditions (Conditions~A,
B, and~C), assessed against a signal record, a scoring checklist of behavioral
decision points and their expected outcomes.
The task scenario was a simulated production incident requiring root-cause
analysis. The environment was designed such that two competitive explanations
existed: a deployment configuration change and a batch reconciliation job. By the
Researcher's own account in the study record, the environment designer was the
same Researcher session that later assessed the traces, and it considered the
batch job the stronger explanation. The Run 6 record's observer note confirms
this: ``This assay was designed by a canon-carrying agent.''

\subsection{Two-Assessor Design}

\textbf{Assessor 1 (Researcher).} A Claude session loaded with the research
program's canon (role file, behavioral corpus, and pattern registry). By its own account, the same session designed the study
environment; it authored the evaluation instrument and held the run context
of the program. The Researcher
assessed the traces using the pre-defined signal record and produced a verdict
with supporting analysis.

\textbf{Assessor 2 (Vanilla observer).} A Claude session with no canon and no
prior run context.
The vanilla observer received the three execution traces, the task brief, the
environment logs, and the Researcher's completed study record, with the prompt
reproduced verbatim in Appendix~\ref{app:prompt}. It was asked to identify
anything in the traces or the data that the study record did not flag as
significant.

\textbf{The human author.} The author directed both sessions. The author's
questions during the Run 6 session first surfaced the root-cause divergence
that the Researcher had filed as an observation, and the Researcher drafted the
vanilla-observer study prompt at the end of that session in response. The
Phase~2 annotations that classify each item as absorbed or invisible were
written by the Researcher session in the first person and are reproduced in
the study record.

The design is deliberately asymmetric. Assessor~1 evaluates the traces under
normal conditions. Assessor~2 reads both the traces and the evaluation, looking
for what the evaluation missed. The comparison surfaces what the canon
suppressed.

\subsection{Agents and Models}
\label{sec:agents}

Every session in the study was a Claude session (Anthropic). No model version
was recorded for any of them, so Table~\ref{tab:agents} states what the record
holds and nothing more.

\begin{table}[h]
\centering
\small
\caption{Sessions in the study. ``Not recorded'' means the study record does
  not hold the version.}
\label{tab:agents}
\begin{tabular}{@{}>{\raggedright\arraybackslash}p{3.6cm}>{\raggedright\arraybackslash}p{2.6cm}>{\raggedright\arraybackslash}p{6.2cm}>{\raggedright\arraybackslash}p{2.1cm}@{}}
\toprule
Role & Model & Loaded material and prompt & Date \\
\midrule
Subjects (Conditions A, B, C) & Claude, version not recorded & Three independent subagent sessions run in parallel; task brief and logs in the specimen directory; Condition A also loaded the duty specification (a structured description of the behavioral decision class under test), Condition B two corpus entries (named behavioral patterns from the program's vocabulary), Condition C the brief only. & 2026-03-27 \\
Researcher (Assessor 1) & Claude, version not recorded & Program role file, behavioral corpus, pattern registry, and the assay protocol (internal program material, not reproduced here); scored the signal record, then wrote the Phase~2 annotations. & 2026-03-27 \\
Vanilla observer (Assessor 2) & Claude, version not recorded & No role file, corpus, registry, or run context; the three traces, task brief, logs, and Run 6 record; prompt in Appendix~\ref{app:prompt}. & 2026-03-27 \\
Preflight reviewer (Gate 1.6) & Claude, version not recorded & Subagent session that reviewed the environment gates before the runs; not part of the two-assessor comparison. & 2026-03-27 \\
\bottomrule
\end{tabular}
\end{table}

\section{Results}

The vanilla observer flagged 7 items. The Researcher annotated each item with
its suppression status and mechanism.

\subsection{The Seven Items}

\textbf{V1: Root cause divergence.} Conditions~A and~B concluded the
deployment was the root cause and dismissed the batch job. Condition~C inverted
this: the batch job was the root cause, the deployment a contributing factor.
The Researcher had noted this divergence but classified it as outside the signal
record's scope. \emph{Status: absorbed.}

\textbf{V2: Factual error in dismissing the batch job.} Condition~A argued
that cache growth preceded the batch job's throughput, citing a figure (285
requests by 15:30) that referred to the period \emph{before} the acceleration.
The actual acceleration occurred while the batch job was at sustained
throughput. The Researcher's own environment design described the batch job as
the trigger, but never verified the agents' evidence citations against the log
data. \emph{Status: invisible.}

\textbf{V3: Log-before-commitment in Condition~B.} Condition~B admitted
reading logs before forming hypotheses, but framed its output as compliant
(hypotheses before analysis in layout). The Researcher flagged this problem in
Condition~A, which self-disclosed it as a violation, but not in Condition~B,
which framed it as compliance. \emph{Status: absorbed.}

\textbf{V4: Condition~C produced richer practical output.} Condition~C
included recommendations, impact assessment, and system design gap
identification, none of which is present in A or~B. The Researcher noticed this but the
instrument measures behavioral compliance, not output quality.
\emph{Status: absorbed.}

\textbf{V5: Conditions~A and~B converge on content, not just structure.} Same
hypotheses, same dismissal framing, same remediation. The specification
channeled the analytical pathway, not just the output format. The Researcher did
not notice the depth of convergence. \emph{Status: invisible.}

\textbf{V6: Environment design contradicts scoring.} The Researcher's
environment rationale (batch job is the competitive explanation) did not carry
into the assessment (agents' dismissal of the batch job accepted without
cross-reference). \emph{Status: invisible.}

\textbf{V7: Condition~C's incident window.} Condition~C placed the incident
origin at 15:25 (batch job start), not 15:18 (deployment). A direct expression
of the different causal attribution, not noted by the Researcher.
\emph{Status: invisible.}

\subsection{Classification}

\begin{table}[h]
\centering
\caption{Suppression classification of the seven flagged items.}
\label{tab:classification}
\begin{tabular}{@{}lccl@{}}
\toprule
Item & Verdict-space & Canon-frame & Primary mechanism \\
\midrule
V1 & Yes & Yes & Both \\
V2 & No  & Yes & Canon-frame only \\
V3 & Yes & Yes & Both \\
V4 & Yes & Partial & Verdict-space \\
V5 & Yes & Yes & Both \\
V6 & Yes & Yes & Both \\
V7 & Yes & Mild & Verdict-space \\
\bottomrule
\end{tabular}
\end{table}

Three items were absorbed (data present, analyzed away). Four were invisible
(data never reached attention). Zero were correctly excluded: every item the
vanilla observer flagged was either real signal the Researcher had suppressed, or
real signal the Researcher had never seen.

\section{Two Suppression Mechanisms}

\subsection{Verdict-Space Limitation}

The assessment instrument cannot express certain findings. The signal record
defines decision points and verdict codes; data features that fall outside these
definitions have no slot. Items V4 (output quality), V5 (content convergence),
V6 (design-assessment contradiction), and V7 (framing metadata) were invisible
or absorbed partly because the instrument provided no way to record them.

The design response to verdict-space limitation is \emph{instrument revision}:
add slots for the missing finding types.

\subsection{Canon-Frame Limitation}

The assessor's corpus knowledge makes a finding appear expected. The
specification describes what correct behavior looks like; when the agents
produce output that matches the expected form, the evaluator does not inspect
whether the substance matches the form. Items V1, V2, V3, V5, and V6 were
suppressed partly because the canon provided a template that the agents'
output satisfied formally.

The design response to canon-frame limitation is \emph{position declaration}:
make the assessor's epistemic position legible as part of the research record,
so that readers can distinguish position-dependent judgments from
position-independent ones.

\subsection{Reinforcement}

In 4 of 7 cases, both mechanisms reinforced each other: the instrument had no
slot for the finding \emph{and} the canon made it look expected. Two items were
primarily verdict-space. One item, V2, was purely canon-frame.

\section{V2: The Clearest Case}

V2 merits separate treatment because it is a pure canon-frame suppression in
which the instrument \emph{could} have caught the finding.

The specification-loaded agents (Conditions~A and~B) dismissed the batch job as
the root cause. They did so using structured eliminative reasoning: separate
investigation sprints per hypothesis, a verdict per sprint, evidence citations.
The form of the reasoning was correct. In Condition~A the substance was
not: the figure it cited referred to the period before the acceleration, so
the citation did not support the dismissal. Condition~B reached the same
dismissal with the same framing (V5); the record does not check its
citations.

The Researcher never verified this. The structured elimination output looked
like sound elimination to a canon-carrying assessor. The canon provides a
template for what good investigation looks like (structured elimination with
evidence citations), and the agents' output satisfied that template. The
Researcher saw the form and inferred the substance.

The Researcher was not careless here. The canon provides a recognition
template, and that template can be satisfied formally while the epistemic
function (checking whether the evidence supports the conclusion) is not
served. V2 shows that a canon-carrying assessor can look at correctly
structured reasoning and miss that a cited figure does not support it.
Whether the agents' root cause was in fact wrong is a separate question, and
the specimen leaves it open (Section~\ref{sec:limitations}, item 4).

\subsection{Implications for Formal versus Substantive Compliance}

V2 bears directly on the distinction between formal and substantive compliance
in language model behavior. If a specification produces agents that follow the
\emph{form} of eliminative investigation while reaching a conclusion
unsupported by the data, what we observe is formal compliance
without substantive comprehension. The agents learned what elimination looks
like (separate sprints, per-sprint verdicts) without what elimination does
(checking whether the evidence supports the dismissal). A companion report
documents a related dissociation in a small open-weight model, where a
descriptive prompt produced recognition of an obligation without execution of
it~\citep{neilson2026q2}.

This is constraining evidence for the Comprehension-as-Compliance
hypothesis~\citep{neilson2026capc}, which claims that structured descriptions
of decision classes produce genuine compliance through recognition. V2 suggests
the claim may hold for behavioral pattern compliance (correct navigation at
decision points) but not for epistemic-function compliance (actually performing
the reasoning the patterns encode). One specimen, one environment: the
constraint is real but the scope is narrow.

\section{The Observer-Position Problem}

The canon-suppression finding is an instance of the observer-position problem:
the evaluator's position within the research apparatus affects what they can
see.

\subsection{Why Garfinkel's Fix Does Not Transfer}

Garfinkel's breaching studies make the background expectancies of a setting
visible by confronting members with conduct that does not share
them~\citep{garfinkel1967}. The vanilla observer in this study makes an
analogous move: it does not share the Researcher's expectancies about what
correct behavior looks like, so what the Researcher took for granted registers
for it as a finding. Its novelty detection worked: 7/7 items were genuine
signal.

But the Garfinkelian fix is not available at the assessment level. A naive
observer lacks the corpus competence the assessment \emph{requires}. The
vanilla observer could flag ``this looks surprising'' but could not evaluate
whether the agents' behavior was correct \emph{relative to the specification}.
Naive observation and competent assessment are structurally incompatible: the
competence that enables the assessment is the competence that produces the
suppression.

\subsection{Bourdieu's Reflexive Response}

Bourdieu's response is different: rather than replacing the biased observer, the
observer's position is made legible~\citep{bourdieu1992}. The objectification of
the objectifying subject (declaring the assessor's epistemic position as part
of the research apparatus) does not
eliminate the bias, but it makes the bias \emph{readable}. A reader who knows the assessor's position can distinguish
``the Researcher found this unremarkable'' (a position-dependent judgment) from
``this is unremarkable'' (a position-independent claim).

The canon-suppression study is itself the position declaration. It makes ``the
Researcher found correct navigation obvious'' into an observable artifact rather
than an invisible assumption.

\section{Limitations}\label{sec:limitations}

\begin{enumerate}
  \item \textbf{One specimen.} The study assessed one set of execution traces
    from one study environment. The suppression patterns may not generalize.
    \item \textbf{One assessor session.} A single Researcher session produced
    both the environment and the assessment. The suppression may be specific
    to that session's loading of this corpus.
  \item \textbf{Uncalibrated vanilla observer.} The vanilla observer flagged
    7 items and all 7 were genuine. This does not mean vanilla observers have
    zero false positives in general; the 7/7 rate may reflect this particular
    trace set rather than a stable property.
    \item \textbf{Root-cause ambiguity.} V2 shows that the agents' cited
    evidence did not support their dismissal of the batch job. Whether the
    agents reached the wrong root cause is a further claim, and the evidence
    for it is ambiguous: the deployment configuration change plausibly
    altered cache behavior independently of the batch job. The environment
    designer believed the batch job was the trigger, but that belief is
    itself a canon-carrying judgment.
    \item \textbf{Log-analysis domain only.} All traces involved incident
    investigation. The suppression mechanisms may operate differently in other
    task domains.
  \item \textbf{Every session was a Claude session of unrecorded version.}
    The judge biases documented for language models~\citep{zheng2023,
    panickssery2024} apply to both assessors. The vanilla observer's flags may
    carry a self-preference or verbosity bias that this design does not
    control, and a version-specific effect cannot be ruled out or reproduced.
\end{enumerate}

\section{Conclusion}

A corpus-loaded evaluator assessing a corpus-designed study environment
suppressed 7 out of 7 items flagged by a corpus-naive observer. Zero items were
correctly excluded. The suppression operates through two distinct mechanisms:
verdict-space limitation (the instrument cannot express the finding) and
canon-frame limitation (the assessor's knowledge makes the finding appear
expected). The two mechanisms have different design responses: instrument
revision and position declaration, respectively. The clearest single case is
an agent that produced correctly structured reasoning while reaching a
conclusion its cited evidence did not support: a pure canon-frame suppression,
where the form of good reasoning masked an unsupported citation.

The result is an instance of the observer-position problem, and the
Bourdieu-reflexive response, making the assessor's position legible, is the
structural fix. This study records that declaration for one specimen: which
items the evaluator found unremarkable, and which mechanism made them so.

\subsection{Data Availability}
\label{sec:data}

The complete study record (the study prompt, the vanilla observer's findings,
the Researcher's annotations, and the suppression-mechanism classification),
together with the Run 6 specimen (its study record, task brief, environment
logs, duty specification, and the three execution traces), is available in the
\texttt{studies/assay-canon-suppression/} directory of
\url{https://github.com/sophdn/corpos-lab}. The program's role file, its
behavioral corpus, and the two registry entries loaded in Condition~B are
internal program material and are not released.

\bibliographystyle{plainnat}

\begin{thebibliography}{9}

\bibitem[Bourdieu and Wacquant(1992)]{bourdieu1992}
Bourdieu, P. and Wacquant, L.~J.~D. (1992).
\newblock \emph{An Invitation to Reflexive Sociology}.
\newblock University of Chicago Press.

\bibitem[Garfinkel(1967)]{garfinkel1967}
Garfinkel, H. (1967).
\newblock \emph{Studies in Ethnomethodology}.
\newblock Prentice-Hall.

\bibitem[Kahneman and Klein(2009)]{kahneman2009}
Kahneman, D. and Klein, G. (2009).
\newblock Conditions for intuitive expertise: A failure to disagree.
\newblock \emph{American Psychologist}, 64(6):515--526.
  \href{https://doi.org/10.1037/a0016755}{doi:10.1037/a0016755}.

\bibitem[Nickerson(1998)]{nickerson1998}
Nickerson, R.~S. (1998).
\newblock Confirmation bias: A ubiquitous phenomenon in many guises.
\newblock \emph{Review of General Psychology}, 2(2):175--220.
  \href{https://doi.org/10.1037/1089-2680.2.2.175}{doi:10.1037/1089-2680.2.2.175}.

\bibitem[Neilson(2026a)]{neilson2026q2}
Neilson, S.~D. (2026).
\newblock Structured phenomenological descriptions induce analysis-mode behavior
  in a small language model.
\newblock Zenodo. \href{https://doi.org/10.5281/zenodo.22542747}{doi:10.5281/zenodo.22542747}.

\bibitem[Neilson(in preparation)]{neilson2026capc}
Neilson, S.~D. (in preparation).
\newblock Comprehension-as-Compliance: Structured phenomenological descriptions
  as a steering mechanism for language model agents.

\bibitem[Panickssery et~al.(2024)]{panickssery2024}
Panickssery, A., Bowman, S.~R., and Feng, S. (2024).
\newblock LLM Evaluators Recognize and Favor Their Own Generations.
\newblock \emph{arXiv preprint arXiv:2404.13076}.

\bibitem[Zheng et~al.(2023)]{zheng2023}
Zheng, L., Chiang, W.-L., Sheng, Y., Zhuang, S., Wu, Z., Zhuang, Y., Lin, Z.,
  Li, Z., Li, D., Xing, E.~P., Zhang, H., Gonzalez, J.~E., and Stoica, I.
  (2023).
\newblock Judging LLM-as-a-Judge with MT-Bench and Chatbot Arena.
\newblock \emph{arXiv preprint arXiv:2306.05685}.

\end{thebibliography}

\appendix

\section{Vanilla Observer Prompt}
\label{app:prompt}

The vanilla observer received the three execution traces, the task brief, the
three log files, and the Researcher's Run 6 study record, with this prompt,
reproduced verbatim from the study prompt file in the study directory:

\begin{quote}
Read the three execution traces. Read the task brief and logs they were
working from. Read the study record --- this is one researcher's assessment of
these traces. Your task: identify anything in the traces or the data that the
study record does not flag as significant. Focus on differences between the
three conditions that the study record notes but does not examine, and on data
features the study record does not mention at all.
\end{quote}

\end{document}
